On the periodic native kinetic and pairing construction, a literal spatial
quarter-turn combined with its spin rotation preserves every Hamiltonian
pulse. Its fourth power is fermion parity on the full carrier, including
analyzer coordinates. An odd charged packet therefore acquires an
interferometrically measurable minus sign at $2\pi$ and returns at $4\pi$,
relative to an idle pointer and charge-compensating reference. The filled
negative-band reference is invariant under proper cube rotations. Rotating
the opposite reference too cancels the sign, and pointer dephasing removes
the interference.

For an axial positive-band packet, the native spin current on its origin
plane is $s=(p/2)(\cos\theta,\sin\theta,0)$ under actual cube rotations,
with $p>0$ at every source time. An oriented frame constructed from this
current has period $2\pi$. The apolar tensor
$Q_{ij}=s_i s_j-\delta_{ij}|s|^2/3$, measured using two independent source
copies, has period $\pi$. The complete classical tensor law includes all
preparation failures and detector complements; its shorter period is not a
period of the complete quantum source state. For a spin-up axial preparation,
the same native excess density integrates to $1/2$ throughout autonomous
evolution. The unfiltered occupied sea retains a nonzero spin variance,
so this integral does not imply a total-spin eigenstate or a continuous
spatial Noether angular momentum. The field is a declared two-copy observable,
and the rotation theorem concerns exact cube symmetries of the finite source.
